\documentclass[12pt,a4paper]{article}
\usepackage{fontspec}
\usepackage{polyglossia}
\setdefaultlanguage{russian}
\setotherlanguage{english}
\setmainfont{STIX Two Text}
\newfontfamily\cyrillicfonttt{Menlo}
\usepackage[a4paper,margin=2cm]{geometry}
\usepackage{graphicx}
\usepackage{float}
\usepackage{booktabs}
\usepackage{longtable}
\usepackage{tabularx}
\usepackage{array}
\usepackage{enumitem}
\usepackage{hyperref}
\usepackage{caption}
\usepackage{pdflscape}
\usepackage{xurl}
\hypersetup{colorlinks=true,linkcolor=blue,urlcolor=blue}
\setlength{\parindent}{1.25cm}
\setlength{\parskip}{0.25em}
\renewcommand{\arraystretch}{1.18}
\setlist{nosep,leftmargin=1.5em}
\graphicspath{{report-assets/lab2/}}
\newcolumntype{Y}{>{\raggedright\arraybackslash}X}
\newcommand{\shot}[2]{%
\begin{figure}[H]\centering
\includegraphics[width=0.96\textwidth,height=0.36\textheight,keepaspectratio]{report-assets/lab2/#1}%
\caption{#2}\end{figure}}
\newcommand{\field}[2]{\noindent\textbf{#1.} #2\par}
\newcommand{\scenario}[7]{%
\subsection*{#1 — #2}
\field{Исходные условия}{#3}
\field{Тестовые данные}{#4}
\field{Действия}{#5}
\field{Ожидаемый ответ}{#6}
\field{Ожидаемое состояние данных}{#7}
\field{Статус}{Сформулировано, не выполнено.}
}
\begin{document}
\begin{titlepage}
\centering
{\large Университетская лабораторная работа\par}
\vspace{2.5cm}
{\LARGE\bfseries Отчет по лабораторной работе № 2\par}
\vspace{0.6cm}
{\Large Формирование требований безопасности к программной системе\par}
\vspace{1.8cm}
\begin{tabular}{p{4cm}p{11cm}}
Дисциплина: & Основы разработки защищенного программного обеспечения и компьютерных сетей\\
Вариант: & 19 — Система учета транспортных средств организации\\
Исходная версия: & v0.1\\
Версия ЛР: & v0.2\\
Студент: & Мостовая Полина Андреевна\\
Группа: & ИУСбд-02-23
\end{tabular}
\vfill
{\large 2026\par}
\end{titlepage}
\tableofcontents
\clearpage

\section{Цель работы}
Цель работы — определить, какие данные и функции системы нужно защищать, кому они доступны и как проверять требования безопасности. Эта лабораторная продолжает проект, начатый в работе № 1. Основной результат — проектная документация; механизмы защиты будут разрабатываться позднее \cite{method2}.

\section{Назначение и границы системы}
Система учета транспортных средств организации предназначена для ведения реестра автопарка, закрепления транспорта за водителями, планирования поездок, учета пробега и контроля технического состояния. Она нужна, чтобы сведения об использовании служебного транспорта хранились и обновлялись в одном месте.

В границы проекта входят приложение, база данных, конфигурация и журнал событий. Кадровая система может быть источником профилей, но интеграция с ней в проект не входит. Система не управляет движением транспорта и не определяет его местоположение. Для учебных проверок используются искусственные данные.

Сейчас приложение загружает и проверяет конфигурацию, затем выводит стартовое сообщение. Интерфейса, базы данных и учета транспорта и поездок пока нет; вход пользователей, разграничение доступа и журналирование также не реализованы.

На снимке показан фактический результат запуска приложения.
\shot{screen10.jpg}{Запуск приложения и стартовое сообщение}

\subsection{Роли}
\begin{description}
\item[Водитель] Просматривает только закрепленные транспортные средства и собственные назначенные поездки, передает пробег и сведения о неисправностях.
\item[Диспетчер] Ведет карточки транспорта, назначает транспорт, планирует поездки, учитывает техническое состояние и обслуживание.
\item[Руководитель] Просматривает агрегированные учетные отчеты без деталей маршрутов.
\item[Администратор] Управляет учетными записями, ролями и настройками, просматривает журнал безопасности. Автоматического доступа ко всем поездкам у него нет.
\item[Неаутентифицированный пользователь] Может обратиться только к форме входа; данные и остальные операции ему недоступны.
\end{description}

\subsection{Функции}
Функции F-1--F-8 являются проектируемыми: вход, просмотр и изменение карточек, назначения, поездки, учет технического состояния, отчеты и журнал событий. Загрузка и проверка конфигурации (F-9), а также вывод стартового сообщения (F-10) уже реализованы.

\small
\begin{longtable}{p{0.8cm}p{3.6cm}p{2.7cm}p{4cm}p{3cm}}
\toprule ID & Функция & Пользователь & Основные данные & Состояние \\
\midrule\endhead
F-1 & Вход и проверка учетных данных & Все роли & Идентификатор входа, данные проверки & Запланирована \\
F-2 & Просмотр карточек транспорта & Водитель, диспетчер & Госномер, модель, статус & Запланирована \\
F-3 & Создание и изменение карточек ТС & Диспетчер & Карточка ТС, статус & Запланирована \\
F-4 & Назначение ТС и планирование поездок & Диспетчер & Водитель, ТС, дата, маршрут & Запланирована \\
F-5 & Регистрация поездки, пробега и неисправности & Водитель, диспетчер & Поездка, пробег, результат & Запланирована \\
F-6 & Учет состояния и обслуживания & Диспетчер & Статус ТС, дефект, работы, срок ТО & Запланирована \\
F-7 & Формирование учетных отчетов & Диспетчер; руководитель читает сводку & Автопарк, поездки, обслуживание & Запланирована \\
F-8 & Журналирование действий и событий безопасности & Администратор & Субъект, действие, время, результат & Запланирована \\
F-9 & Загрузка и проверка конфигурации & Администратор среды & APP\_NAME, APP\_ENV, APP\_SECRET\_KEY & Реализована \\
F-10 & Вывод стартового сообщения без секрета & Запускающий приложение & Название, режим работы & Реализована \\
\bottomrule
\end{longtable}
\normalsize

\subsection{Категории обрабатываемых данных}
\small
\begin{longtable}{p{0.8cm}p{2.7cm}p{3.5cm}p{3cm}p{4.2cm}}
\toprule ID & Категория & Примеры полей & Источник и хранение & Ограничение доступа \\
\midrule\endhead
DT-1 & Профили пользователей & Имя, табельный номер, контакт & Администратор; база данных (планируется) & Собственный профиль; администратору только нужные поля \\
DT-2 & Учетные данные и роли & Идентификатор, хэш пароля, роль & Пользователь или администратор; база данных (планируется) & Проверочные данные не отображаются; роль меняет администратор \\
DT-3 & Конфигурация & Режим, секретный параметр & Администратор среды; окружение или локальный .env & Только обслуживающий субъект; секрет не выводится \\
DT-4 & Карточки транспорта & Госномер, модель, статус & Диспетчер; база данных (планируется) & Диспетчер по автопарку; водитель по закреплению \\
DT-5 & Назначения и поездки & Водитель, ТС, маршрут, пробег & Диспетчер и водитель; база данных (планируется) & Диспетчер по задаче; водитель только по своим поездкам \\
DT-6 & Техническое состояние & Дефект, работы, сроки ТО & Водитель сообщает, диспетчер учитывает; база данных (планируется) & Диспетчер изменяет запись; водитель сообщает о дефекте \\
DT-7 & Журналы и отчеты & Действие, время, результат, сводные показатели & Приложение и диспетчер; журнал и база данных (планируются) & Журнал — администратору; руководителю — агрегированные отчеты \\
\bottomrule
\end{longtable}
\normalsize

\section{Активы и их значимость}
Оценки даны по шкале 1--3: 1 — низкая, 2 — средняя, 3 — высокая значимость. $C_i$ — конфиденциальность, $I_i$ — целостность, $A_i$ — доступность. $K_i=\max(C_i,I_i,A_i)$ показывает наибольшую оценку, но не является оценкой риска.

\begin{landscape}
\small
\begin{longtable}{p{1.1cm}p{3.2cm}p{3.2cm}ccccp{9.2cm}}
\toprule ID & Актив & Ответственный & $C_i$ & $I_i$ & $A_i$ & $K_i$ & Обоснование оценок \\
\midrule\endhead
AS-1 & Учетные записи и данные проверки пароля & Администратор безопасности & 3 & 3 & 3 & 3 & Раскрытие может открыть учетные записи; подмена нарушает идентификацию; без входа пользователи не работают. \\
AS-2 & Профили и роли & Администратор безопасности & 2 & 3 & 2 & 3 & Раскрываются служебные сведения; неверная роль меняет полномочия; восстановление профилей требует ручной работы. \\
AS-3 & Секрет и конфигурация приложения & Администратор среды & 3 & 3 & 2 & 3 & Раскрытие ослабляет защиту; подмена меняет режим; конфигурацию можно восстановить из защищенного источника. \\
AS-4 & Реестр транспорта & Владелец процесса учета автопарка & 2 & 3 & 3 & 3 & Искажение карточек нарушает учет; недоступность останавливает основную функцию. \\
AS-5 & Назначения и поездки & Владелец диспетчеризации & 3 & 3 & 3 & 3 & Маршруты связаны с сотрудниками; подмена нарушает учет; потеря доступа блокирует диспетчеризацию. \\
AS-6 & Техническое состояние и обслуживание & Ответственный за технический учет & 2 & 3 & 3 & 3 & Неверный статус влияет на эксплуатационные решения; недоступность мешает контролю исправности. \\
AS-7 & Журнал событий безопасности & Ответственный за контроль безопасности & 2 & 3 & 2 & 3 & Подмена мешает разбору событий; краткая недоступность допускает отложенный анализ. \\
AS-8 & Учетные отчеты & Владелец процесса учета автопарка & 2 & 3 & 2 & 3 & Искажение влияет на решения; отчет можно восстановить из первичных данных. \\
AS-9 & Исходный код и тесты & Разработчик проекта & 1 & 3 & 2 & 3 & Подмена может внести дефекты; исходники восстанавливаются из Git. \\
AS-10 & Работоспособность системы & Владелец сервиса & 1 & 2 & 3 & 3 & Длительная остановка блокирует основные операции; краткий сбой не искажает сохраненные данные. \\
\bottomrule
\end{longtable}
Активы пока проектируются. Сейчас приложение только загружает конфигурацию и проверяет длину значения секретного ключа; криптографически ключ не используется.
\end{landscape}

\begin{landscape}
\section{Предварительная матрица доступа}
Неаутентифицированный пользователь может открыть форму входа, но не получает данные. Явно не разрешенные действия запрещены. Роли не наследуются. Для водителя проверяется принадлежность поездки/транспорта; руководитель видит только агрегированные отчеты. Администратор не получает автоматического доступа к маршрутам. Эти правила соответствуют рекомендациям OWASP проверять полномочия для каждого запроса и придерживаться минимальных привилегий \cite{owasp-auth}.

\small
\begin{tabularx}{\linewidth}{p{6cm}YYYYY}
\toprule Операция & Без входа & Водитель & Диспетчер & Руководитель & Администратор \\
\midrule
Форма входа & Разрешено, без выдачи данных & Разрешено & Разрешено & Разрешено & Разрешено \\
Просмотр карточки ТС & Запрещено & Только закрепленное ТС & Весь автопарк по задаче & Запрещено & Запрещено, кроме минимальных полей для учетной записи \\
Создать/изменить карточку ТС & Запрещено & Запрещено & Разрешено & Запрещено & Запрещено \\
Просмотр маршрута & Запрещено & Только собственная назначенная поездка & По автопарку в рамках работы & Запрещено & Запрещено по умолчанию \\
Изменение назначения & Запрещено & Запрещено & Разрешено & Запрещено & Запрещено \\
Передача отчета о неисправности & Запрещено & Только по своей поездке & В рамках учета & Запрещено & Запрещено \\
Изменение обслуживания & Запрещено & Запрещено; можно сообщить о дефекте & Разрешено & Запрещено & Запрещено \\
Просмотр учетного отчета & Запрещено & Запрещено & По служебной задаче & Только агрегированные показатели & Запрещено по умолчанию \\
Учетные записи и роли & Запрещено & Запрещено & Запрещено & Запрещено & Разрешено по процедуре \\
Просмотр журнала безопасности & Запрещено & Запрещено & Запрещено & Запрещено & Разрешено в рамках обязанностей безопасности \\
\bottomrule
\end{tabularx}
Матрица описывает планируемую политику и пока не реализована в коде.
\end{landscape}

\section{Требования безопасности}
Подготовлено 13 требований: восемь функциональных, три нефункциональных и два процессных. Приоритет для каждого — обязательный. Далее приведены критерии, способы проверки и текущие статусы.

\subsection*{SEC-1. Доступ после аутентификации}
\field{Тип}{Функциональное.}
\field{Активы}{AS-1, AS-2, AS-4--AS-8.}
\field{Обоснование}{Не допустить анонимный доступ к сведениям о сотрудниках, транспорте и поездках.}
\field{Критерий}{Анонимный запрос защищенной функции отклонен; содержимое не возвращается, данные не меняются.}
\field{Проверка}{Негативный интеграционный сценарий AT-1 и тест успешного входа.}
\field{Статус}{Сформулировано.}

\subsection*{SEC-2. Проверка полномочий по роли}
\field{Тип}{Функциональное.}
\field{Активы}{AS-2, AS-4--AS-8.}
\field{Обоснование}{Предотвратить действия вне служебных полномочий.}
\field{Критерий}{Разрешенная операция выполнена; запрещенная отклонена без изменения данных.}
\field{Проверка}{Тесты ролей по матрице доступа; сценарии AT-3--AT-6.}
\field{Статус}{Сформулировано.}

\subsection*{SEC-3. Доступ водителя к назначенным данным}
\field{Тип}{Функциональное.}
\field{Активы}{AS-2, AS-4, AS-5.}
\field{Обоснование}{Ограничить раскрытие чужих поездок и назначений.}
\field{Критерий}{Собственная поездка доступна; запрос чужой поездки отклоняется без выдачи полей и изменения данных.}
\field{Проверка}{Пара AT-3/AT-4 с двумя водителями.}
\field{Статус}{Сформулировано.}

\subsection*{SEC-4. Изменение карточки транспортного средства}
\field{Тип}{Функциональное.}
\field{Активы}{AS-4.}
\field{Обоснование}{Защитить целостность реестра.}
\field{Критерий}{Диспетчер сохраняет изменение; водитель и администратор получают отказ, карточка остается прежней.}
\field{Проверка}{Тест ролей и проверка записи до и после запроса.}
\field{Статус}{Сформулировано.}

\subsection*{SEC-5. Изменение назначения и плана поездки}
\field{Тип}{Функциональное.}
\field{Активы}{AS-2, AS-5.}
\field{Обоснование}{Сохранить правильность планирования и закрепления транспорта.}
\field{Критерий}{Изменение диспетчера сохраняется; запрос другой роли отклонен, назначение не меняется.}
\field{Проверка}{Интеграционные тесты ролей и состояния записи.}
\field{Статус}{Сформулировано.}

\subsection*{SEC-6. Проверка отчета о поездке}
\field{Тип}{Функциональное.}
\field{Активы}{AS-2, AS-5.}
\field{Обоснование}{Ограничить подмену чужих отчетов и ошибочные значения пробега.}
\field{Критерий}{Корректный пробег своей поездки сохраняется; отрицательный пробег или чужая поездка отклоняются без создания/изменения отчета.}
\field{Проверка}{Позитивный и негативный сценарии AT-5 и AT-6.}
\field{Статус}{Сформулировано.}

\subsection*{SEC-7. Журналирование значимых действий}
\field{Тип}{Функциональное.}
\field{Активы}{AS-1, AS-2, AS-5, AS-7.}
\field{Обоснование}{Сохранить сведения для последующего разбора действий.}
\field{Критерий}{Для тестового изменения и отказа созданы записи с субъектом, действием, временем и результатом.}
\field{Проверка}{Проверить журнал после успешной и отклоненной операции.}
\field{Статус}{Сформулировано.}

\subsection*{SEC-8. Доступ к журналу безопасности}
\field{Тип}{Функциональное.}
\field{Активы}{AS-2, AS-7.}
\field{Обоснование}{Ограничить раскрытие истории действий сотрудников.}
\field{Критерий}{Администратор безопасности получает журнал; водитель и диспетчер не видят его содержимое.}
\field{Проверка}{Тест доступа к журналу для каждой роли.}
\field{Статус}{Сформулировано.}

\subsection*{SEC-9. Отсутствие секрета в выводе}
\field{Тип}{Нефункциональное.}
\field{Активы}{AS-3, AS-9.}
\field{Обоснование}{Не раскрывать секрет через консольные сообщения.}
\field{Критерий}{При штатном запуске и предусмотренных ошибках тестовое значение секрета отсутствует в stdout/stderr.}
\field{Проверка}{Захват stdout/stderr с искусственным секретом.}
\field{Статус}{Реализовано частично: стартовое сообщение не включает секрет; штатный вывод покрыт тестами, вывод ошибок отдельно не проверялся.}

\subsection*{SEC-10. Проверка конфигурации до запуска}
\field{Тип}{Нефункциональное.}
\field{Активы}{AS-3, AS-10.}
\field{Обоснование}{Не допустить запуск в неопределенном режиме.}
\field{Критерий}{Пустое имя, неизвестная среда или ключ короче 32 символов приводят к контролируемой ошибке до старта.}
\field{Проверка}{Тесты конфигурации для трех ошибочных значений; анализ кода.}
\field{Статус}{Реализовано: проверки есть в app/config.py; отдельные автоматические тесты не добавлены.}

\subsection*{SEC-11. Защита сетевой передачи}
\field{Тип}{Нефункциональное.}
\field{Активы}{AS-1, AS-5, AS-10.}
\field{Обоснование}{Снизить риск перехвата учетных данных и поездок.}
\field{Критерий}{После появления сетевого интерфейса тестовое развертывание использует TLS и не принимает учетные данные по незащищенному каналу.}
\field{Проверка}{Проверка конфигурации развертывания и сетевое испытание.}
\field{Статус}{Сформулировано; сетевого интерфейса пока нет.}

\subsection*{SEC-12. Проверка изменений до выпуска}
\field{Тип}{Процессное.}
\field{Активы}{AS-9, AS-1--AS-8.}
\field{Обоснование}{Снизить вероятность внесения дефектов в проект и защитные механизмы.}
\field{Критерий}{Для изменения имеются результат ревью и результаты применимых тестов до включения в релиз.}
\field{Проверка}{Анализ истории изменений и протокола проверок.}
\field{Статус}{Сформулировано.}

\subsection*{SEC-13. Работа с секретами в Git}
\field{Тип}{Процессное.}
\field{Активы}{AS-3, AS-9.}
\field{Обоснование}{Не допустить публикацию реальных секретов.}
\field{Критерий}{Файл .env не отслеживается, staged diff не содержит реальных секретов, шаблон содержит условное значение.}
\field{Проверка}{git ls-files .env; проверка .env.example и staged diff.}
\field{Статус}{Реализовано частично: файл исключен, шаблон существует; diff требуется проверять перед фиксацией.}

\section{Сценарии приемки}
Сценарии описывают ожидаемое поведение функций, которые пока не реализованы. Ни один сценарий не выполнялся. Их критерии соотносятся с проверяемыми требованиями OWASP ASVS \cite{owasp-asvs}.

\scenario{AT-1}{Негативный: анонимный запрос карточки ТС}{Система запущена; пользователь не вошел; существует карточка ТС-001.}{Запрос карточки ТС-001 без учетных данных.}{Открыть адрес карточки, записать ответ, затем проверить карточку уполномоченной тестовой учетной записью.}{Запрос отклонен или перенаправлен на вход; поля карточки не раскрываются.}{Карточка не меняется; новые поездки и назначения не создаются.}
\scenario{AT-2}{Позитивный: успешный вход}{Учетная запись DRV-001 активна.}{Искусственные действительные учетные данные DRV-001.}{Отправить данные форме входа; открыть профиль водителя.}{Вход принят; доступны только данные, разрешенные роли.}{Роль и учетные данные не меняются; секрет проверки клиенту не возвращается.}
\scenario{AT-3}{Позитивный: чтение назначенной поездки}{DRV-001 вошел; TRIP-001 назначена ему и связана с ТС-001.}{DRV-001, TRIP-001, ТС-001.}{Водитель запрашивает TRIP-001 и сверяет показанные дату и статус.}{Сведения собственной поездки возвращены.}{Чтение не меняет поездку или назначение.}
\scenario{AT-4}{Негативный: запрос чужой поездки}{DRV-001 и DRV-002 существуют; TRIP-002 назначена DRV-002, статус «Запланирована».}{DRV-001 запрашивает TRIP-002.}{Войти как DRV-001; запросить TRIP-002; проверить запись от диспетчера.}{Отказ без раскрытия маршрута, пробега и деталей.}{Статус и назначенный водитель не изменились.}
\scenario{AT-5}{Позитивный: допустимый пробег}{DRV-001 вошел; TRIP-001 назначена ему; начальный пробег 12400 км.}{Конечный пробег 12525 км; результат поездки заполнен.}{Открыть свою поездку; внести пробег и результат; отправить отчет; прочитать запись.}{Отчет принят.}{Отчет с пробегом 12525 км связан с TRIP-001; чужие записи не меняются.}
\scenario{AT-6}{Негативный: пробег ниже исходного}{DRV-001 вошел; TRIP-001 назначена ему; начальный пробег 12400 км.}{Конечный пробег 12300 км.}{Отправить отчет; зафиксировать ответ; повторно прочитать поездку и отчет.}{Отчет отклонен с сообщением о недопустимом пробеге.}{Отчет не создан, прежний пробег не изменен.}

\subsection{Прослеживаемость}
\begin{longtable}{p{2.5cm}p{2.3cm}p{2.5cm}p{7cm}}
\toprule Активы & Функция & Требование & Проверка \\
\midrule\endhead
AS-1, AS-2, AS-4--AS-8 & F-1, F-2 & SEC-1 & AT-1, AT-2 \\
AS-2, AS-4, AS-5 & F-2, F-5 & SEC-3 & AT-3, AT-4 \\
AS-2, AS-5 & F-5 & SEC-6 & AT-5, AT-6 \\
AS-2, AS-4--AS-8 & F-1--F-8 & SEC-2 & Тесты ролей по матрице \\
AS-4 & F-3 & SEC-4 & Тест ролей и состояния карточки \\
AS-2, AS-5 & F-4 & SEC-5 & Тест полномочий и назначения \\
AS-1, AS-2, AS-5, AS-7 & F-8 & SEC-7 & Проверка журнала после разрешенной и отклоненной операции \\
AS-2, AS-7 & F-8 & SEC-8 & Тест чтения журнала по ролям \\
AS-3, AS-9 & F-9, F-10 & SEC-9 & Захват вывода с тестовым секретом \\
AS-3, AS-10 & F-9 & SEC-10 & Проверки некорректной конфигурации \\
AS-1, AS-5, AS-10 & F-1, F-5 & SEC-11 & Конфигурация TLS и сетевое испытание после реализации API \\
AS-9, AS-1--AS-8 & F-1--F-8 & SEC-12 & Проверка ревью и результатов тестов \\
AS-3, AS-9 & F-9 & SEC-13 & git ls-files .env, шаблон и staged diff \\
\bottomrule
\end{longtable}

\section{Согласованность и результаты проверок}
В документах используется единая система идентификаторов. Матрица содержит все роли, а права на поездки и транспорт ограничены принадлежностью объекта. Сценарии связаны с требованиями; тестовые результаты для будущих функций не заявляются.

В проекте определены четыре аутентифицированные роли и неаутентифицированный пользователь, десять функций, семь категорий данных, десять активов, десять операций матрицы, тринадцать требований и шесть сценариев (три позитивных и три негативных). У каждого требования указаны все необходимые поля. Оценки активов находятся в диапазоне 1--3; для каждого приведены Ci, Ii, Ai, Ki и обоснование.

Три исходных автоматических теста выполнены успешно: 3 passed. Они проверяют только стартовое сообщение и не подтверждают работу будущих модулей.
\shot{screen1.jpg}{Исходное состояние репозитория и локальные теги}
\shot{screen2.jpg}{Выполнение трех исходных тестов}

\section{Проверка Git}
Основной результат лабораторной зафиксирован тегом v0.2 на commit 4026dd7. После дополнительной сверки документы уточнены в commit 93f0769, для него создан тег v0.2.1; существующий тег v0.2 не перемещался. Commit и оба тега отправлены в удаленный репозиторий. Сейчас origin/main указывает на 93f0769, тег v0.2 — на исходный commit лабораторной, v0.2.1 — на commit с уточнениями. Отчет и снимки хранятся отдельно от проектных файлов.
\shot{screen3.jpg}{Сводка первоначального diff}
\shot{screen4.jpg}{Обновление README с описанием фактического состояния}
\shot{screen5.jpg}{Проверка рабочего дерева и diff}
\shot{screen6.jpg}{Подготовка документов к фиксации}
\shot{screen7.jpg}{Сводка локального commit}
\shot{screen8.jpg}{Создание локального тега v0.2}
\shot{screen9.jpg}{Состояние Git после commit: остались локальные файлы}

\section{Ответы на контрольные вопросы}
\begin{enumerate}
\item Границы показывают, какие компоненты и процессы защищаются проектом, а какие принадлежат внешним системам.
\item Категория группирует похожие данные; актив — объект или ресурс, ценность которого нужно защищать. Актив может объединять несколько категорий.
\item Конфиденциальность нарушается при раскрытии, целостность — при подмене или уничтожении, доступность — когда ресурс нельзя использовать вовремя.
\item Значимость характеризует последствия потери свойства. Риск дополнительно учитывает угрозу, сценарий и вероятность реализации.
\item Максимум скрывает два других значения и различия профилей C/I/A. Поэтому в реестре сохраняются все три оценки и их объяснения.
\item Роль задает набор полномочий; учетная запись относится к конкретному пользователю, которому назначена роль.
\item Роль сама по себе не показывает, какому пользователю принадлежит объект. Водителю разрешены только его поездки.
\item Нет. Администратору предоставляют полномочия, необходимые для его задач; доступ к прикладным данным требует отдельного обоснования.
\item Функциональное требование задает защитное действие системы, нефункциональное — свойство или ограничение ее работы.
\item Они фиксируют правила разработки и сопровождения, например ревью изменений и выполнение тестов до выпуска.
\item Нужны однозначное действие, условия, наблюдаемый критерий и способ проверки.
\item Критерий описывает ожидаемый результат; способ проверки — как получить и оценить этот результат.
\item ID помогает однозначно связывать требование со сценариями, документами и результатами проверки.
\item Функция использует актив; требование задает защитное поведение; проверка сценария подтверждает соответствие критерию.
\item Позитивный сценарий проверяет разрешенное действие; негативный — отказ запрещенного действия или обработку неверных данных.
\item Отказ в ответе не гарантирует, что данные не были изменены. Поэтому после запроса проверяется состояние записи.
\item Длинное значение может быть предсказуемым. Нужны достаточная случайность и безопасные генерация, хранение и применение.
\item «Сформулировано» означает, что требование записано; «реализовано» — механизм существует; «проверено» — имеется результат фактической проверки.
\item Проверяются состав staged файлов, diff, отсутствие секретов и временных файлов, согласованность документации и результаты применимых тестов.
\item Активы, роли, границы и требования помогают определить нарушителей, угрозы, последствия и оценить риски в следующей лабораторной работе.
\end{enumerate}

\section{Вывод}
В работе описаны роли и границы системы учета транспорта, данные и активы, предварительные правила доступа, требования безопасности и сценарии их проверки. Текущий код по-прежнему ограничен запуском приложения и проверкой конфигурации. Три имеющихся теста прошли; сценарии будущих функций не выполнялись. Commit с исправлениями и тег v0.2.1 опубликованы в удаленном репозитории; исходный тег v0.2 сохранен.

\clearpage
\renewcommand{\refname}{Список использованных источников}
\addcontentsline{toc}{section}{Список использованных источников}
\begin{thebibliography}{99}
\bibitem{method2} Методические указания «Лабораторная работа № 2. Формирование требований безопасности к программной системе».
\bibitem{owasp-asvs} \mbox{OWASP} Foundation. Application Security Verification Standard (ASVS), version 5.0.0 [Электронный ресурс]. URL: \url{https://owasp.org/www-project-application-security-verification-standard/} (дата обращения: 15.10.2026).
\bibitem{owasp-auth} \mbox{OWASP} Foundation. Authorization Cheat Sheet [Электронный ресурс]. URL: \url{https://cheatsheetseries.owasp.org/cheatsheets/Authorization_Cheat_Sheet.html} (дата обращения: 15.10.2026).
\end{thebibliography}
\end{document}
